\sectionkicker{Source-backed technical notes}
\subsection*{AutonomyにはGuardrailが要る — Prompt InjectionとData Exfiltration: Technical Notes}
本欄は記事本文で圧縮した一次資料上の情報を、比較・再検証しやすい形で整理したものである。時系列、確認済みの主張、留意点、一次資料URLを掲載する。
\medskip
\subsection*{Theme at a glance}
\addcontentsline{toc}{subsection}{Theme at a glance}
\begin{center}
\footnotesize
\begin{tabularx}{\linewidth}{@{}p{0.20\linewidth}p{0.10\linewidth}p{0.14\linewidth}X@{}}
\toprule
資料 & 位置づけ & 種別 & 時系列 \\
\midrule
Safe URL / agent link data-exfiltration safeguards & 主要資料 & セキュリティ関連 & 2026-01-28 (セキュリティ技術解説) \\
Prompt Injection Attacks on Agentic Coding Assistants & 補足資料 & 論文 & 2026-01-24 (論文投稿) \\
\bottomrule
\end{tabularx}
\end{center}
\normalsize

\begin{technicalnote}{Safe URL / agent link data-exfiltration safeguards}{主要資料}
\begingroup
% reader-facing Technical Notes break policy
\widowpenalty=10000
\clubpenalty=10000
\displaywidowpenalty=10000
\begin{tabularx}{\linewidth}{@{}>{\bfseries}p{0.22\linewidth}X@{}}
組織 & OpenAI \\
種別 & セキュリティ関連 \\
時系列 & 2026-01-28 (セキュリティ技術解説) \\
\end{tabularx}
\smallskip
{\bfseries 一次資料から整理したtechnical points}
\begin{itemize}[leftmargin=1.5em,itemsep=0.35em]
\item \textbf{一次情報で確認できる事実}: OpenAIは1月28日、prompt injectionされた内容がAgentにURL取得を促し、会話由来の秘密情報をURLへ載せて外部へ送出し得るdata-exfiltration脅威を説明した。
\item \textbf{一次情報で確認できる事実}: 説明されたsafeguardは、独立にpublicと確認されたURLのみ自動取得を許し、それ以外はfetchを遮断するか明示的なuser actionを要求する。OpenAIはこれをdefense-in-depthの一層として位置付ける。
\end{itemize}
\begin{minipage}{\linewidth}
% reader-facing Technical Notes generic boundary/source tail group
{\bfseries 読む際の境界}
\begin{itemize}[leftmargin=1.5em,itemsep=0.35em]
\item \textbf{分析上の留意点}: 対象はURL-based exfiltrationであり、任意のweb contentが安全であることやprompt injection全般の解決を示すものではない。
\item \textbf{分析上の留意点}: security propertyはOpenAIによる実装説明であり、本サーベイでは独立したpenetration testを実施していない。
\end{itemize}
\begin{samepage}
{\bfseries 一次資料}
\begingroup\sloppy
\Urlmuskip=0mu plus 2mu\relax
\begin{itemize}[leftmargin=1.5em,itemsep=0.25em]
\item {\scriptsize\url{https://openai.com/index/ai-agent-link-safety}}
\end{itemize}
\endgroup
\end{samepage}
\end{minipage}
\endgroup
\end{technicalnote}
\medskip

\begin{technicalnote}{Prompt Injection Attacks on Agentic Coding Assistants}{補足資料}
\begingroup
% reader-facing Technical Notes break policy
\widowpenalty=10000
\clubpenalty=10000
\displaywidowpenalty=10000
\begin{tabularx}{\linewidth}{@{}>{\bfseries}p{0.22\linewidth}X@{}}
組織 & Narek Maloyan and Dmitry Namiot \\
種別 & 論文 \\
時系列 & 2026-01-24 (論文投稿) \\
\end{tabularx}
\smallskip
{\bfseries 一次資料から整理したtechnical points}
\begin{itemize}[leftmargin=1.5em,itemsep=0.35em]
\item \textbf{一次情報で確認できる事実}: 本論文はagentic coding assistant、skills、tools、protocol ecosystemにまたがるprompt-injection攻撃を整理するSystematization of Knowledgeで、3次元のattack taxonomyを提案する。
\item \textbf{Author claim}: abstractは78研究、42攻撃手法、18防御機構の統合を報告する。攻撃・防御の集計値は著者によるmeta-analysisであり、新規のcontrolled benchmarkではない。
\end{itemize}
\begin{minipage}{\linewidth}
% reader-facing Technical Notes generic boundary/source tail group
{\bfseries 読む際の境界}
\begin{itemize}[leftmargin=1.5em,itemsep=0.35em]
\item \textbf{分析上の留意点}: 異質な研究を束ねたliterature synthesisであるため、集計率を単一の再現可能な実験結果として扱わない。
\item \textbf{分析上の留意点}: 保持sourceは1月時点のarXiv v1であり、その後の研究によって脅威・防御の全体像は更新され得る。
\end{itemize}
\begin{samepage}
{\bfseries 一次資料}
\begingroup\sloppy
\Urlmuskip=0mu plus 2mu\relax
\begin{itemize}[leftmargin=1.5em,itemsep=0.25em]
\item {\scriptsize\url{http://arxiv.org/abs/2601.17548v1}}
\end{itemize}
\endgroup
\end{samepage}
\end{minipage}
\endgroup
\end{technicalnote}
\medskip
